\tzpanchor{0325}
\tzpentryheader{0325}{Gravitomagnetic Coupling Term}

\begingroup\small
\begingroup\scriptsize
\setlength{\tabcolsep}{4pt}
\renewcommand{\arraystretch}{0.92}
\noindent\begin{tabularx}{\linewidth}{@{}p{0.24\linewidth}p{0.36\linewidth}X@{}}
\textbf{UUID} & \multicolumn{2}{l@{}}{\tzpcode{bf705431-926e-5c7b-905a-4b0d03b88975}}\\
\textbf{PiID} & \tzpcode{2092096282925} & \textbf{TZPi}\quad \tzpcode{6}\quad \textbf{PiPE}\quad \tzpcode{332}\\
\textbf{RTEID} & \textbf{Poloidal $\theta$}\quad \tzpcode{1262.0472 rad} & \textbf{Toroidal $\phi$}\quad \tzpcode{02.0795 rad}\\
\textbf{TZPID} & \tzpcode{ID0325} & \textbf{Node Dependencies}\quad \tzpidref{0324}\\
\textbf{Registry} & \tzpcode{registry\_id\_0325} & \textbf{Scale}\quad magnetofluid\\
\textbf{LOC Classification} & QC809 magnetohydrodynamics & QC6 fluid and plasma dynamics\\
\textbf{arXiv Classification} & Primary: physics.plasm-ph & Cross-list: astro-ph.SR, gr-qc\\
\textbf{Primary Support} & Gravitomagnetic Field & Mass Integral\\
\textbf{Closure Support} & Coupling Strength & \\
\end{tabularx}\par
\endgroup
\endgroup

\tzpsection{Canonical Statement}
This term represents the coupling between rotation and gravity, analogous to the magnetic field produced by moving charges.

\tzpsection{Canonical Equation}
\[
\mathbf{B}_{g}=\nabla\times\mathbf{A}_{g}
\]
\[
\mathbf{J}_{m}=\int \rho(\mathbf{r})\,\mathbf{r}\times\mathbf{v}(\mathbf{r})\,d^{3}r
\]

\begin{minipage}[t]{0.49\linewidth}
\tzpsection{Canonical Symbol Key}
\begingroup
\small
\raggedright
\textbf{Source equation symbols.}\par
\begin{itemize}[leftmargin=1.35em,itemsep=1pt,topsep=1pt,parsep=0pt]
    \item $B$: magnetic field magnitude.
    \item $A$: source equation symbol.
    \item $J$: source equation symbol.
    \item $\rho$: density or measure term used by the entry equation.
    \item $v$: velocity, flow speed, or auxiliary comparison variable.
    \item $\mathcal{L}_{g}$: source equation symbol.
\end{itemize}
\endgroup
\end{minipage}\hfill
\begin{minipage}[t]{0.47\linewidth}
\tzpsection{Wolfram Form}
\begingroup
\scriptsize
\raggedright
\textbf{Source Wolfram model.}\par
\begin{Verbatim}[fontsize=\tiny,breaklines=true,breakanywhere=true]
(* TZPID ID0325 generated source Wolfram model *)
ClearAll[K0325, Cr0325, T, R, A, W, I, N];
eq0325$1 = HoldForm[BG == Grad*AG];
eq0325$2 = HoldForm[JM == Integral*rho[r]*r*v[r]*d^3*r];
eq0325$3 = HoldForm[LG == kappaG*JM*BG];

K0325 = HoldForm[{eq0325$1, eq0325$2, eq0325$3}];

N = "normalized TRAWIN closure object for Gravitomagnetic Coupling Term";
I = "Coupling Strength integration/comparison layer";
W = "Mass Integral wave or operator carrier";
A = "Gravitomagnetic Field admissible action/amplitude condition";
R = "Mass Integral rotation/curl preservation layer";
T = "Gravitomagnetic Field local source condition";

Cr0325[expr_] := HoldForm[N[I[W[A[R[T[expr]]]]]]];

Result0325 = Cr0325[K0325];
\end{Verbatim}
\endgroup
\end{minipage}

\par\vspace{0.45\baselineskip}
\noindent\begin{minipage}[t]{\linewidth}
\begingroup
\small
\raggedright
\textbf{TRAWIN Operator Key.}\par
\noindent\textbf{Time/localization operator:} $\mathsf{T}\!\left[\mathrm{local}\right]$ Gravitomagnetic Field local source condition.\par
\noindent\textbf{Rotation/curl operator:} $\mathsf{R}\!\left[\nabla\times\right]$ Mass Integral rotation/curl preservation layer.\par
\noindent\textbf{Action/amplitude operator:} $\mathsf{A}\!\left[\mathrm{admissible}\right]$ Gravitomagnetic Field admissible action/amplitude condition.\par
\noindent\textbf{Wave/d'Alembertian operator:} $\mathsf{W}\!\left[\Box\right]$ Mass Integral wave or operator carrier.\par
\noindent\textbf{Integration operator:} $\mathsf{I}\!\left[\int\right]$ Coupling Strength integration/comparison layer.\par
\noindent\textbf{Normalization operator:} $\mathsf{N}\!\left[\mathrm{normalized}\right]$ normalized TRAWIN closure object for Gravitomagnetic Coupling Term.\par
\endgroup
\end{minipage}

\clearpage
\noindent\textbf{[ID 0325] Gravitomagnetic Coupling Term}\par
\vspace{0.35em}
\tzpsection{Derivation Layer}
\paragraph{Primary Equation.}
\begin{equation}
Q
=
\int_{\Sigma}
J^\mu
\,\mathrm{d}\Sigma_\mu .
\end{equation}

\paragraph{Primary Derivation.}
Let $J^\mu$ be a conserved Trawinistic current satisfying:

\begin{equation}
\nabla^T_\mu J^\mu = 0 .
\end{equation}

The total charge enclosed by a spacelike hypersurface $\Sigma$ is obtained by integrating the current flux through that hypersurface:

\begin{equation}
Q
=
\int_{\Sigma}
J^\mu
\,\mathrm{d}\Sigma_\mu .
\end{equation}

Because the current is conserved, deformation of $\Sigma$ does not change $Q$, provided the hypersurface continues to enclose the same TZP winding structure.

\begin{equation}
\frac{\mathrm{d}Q}{\mathrm{d}t}=0 .
\end{equation}

Thus,

\begin{equation}
\boxed{
Q
=
\int_{\Sigma}
J^\mu
\,\mathrm{d}\Sigma_\mu .
}
\end{equation}

\paragraph{Interpretation.}
The TZP topological charge is conserved because it is the flux of a divergence-free current through a hypersurface enclosing the zero-point structure.

\par\vspace{0.35em}
\noindent\textbf{Derivable Consequences.}\quad
\begingroup
\footnotesize
\raggedright
The canonical equation anchors Gravitomagnetic Coupling Term by connecting the source condition to the derivation layer and proof certificate.
\par\endgroup

\tzpsection{Equation Support}
\noindent\textbf{1. Gravitomagnetic Field}\par
\[
\mathbf{B}_{g}=\nabla\times\mathbf{A}_{g}
\]
\noindent This fixes the gravitomagnetic field induced by the rapidly rotating mass distribution.\par
\noindent\textbf{2. Mass Integral}\par
\[
\mathbf{J}_{m}=\int \rho(\mathbf{r})\,\mathbf{r}\times\mathbf{v}(\mathbf{r})\,d^{3}r
\]
\noindent This identifies the mass-moment integral that weights the field strength.\par
\noindent\textbf{3. Coupling Strength}\par
\[
\mathcal{L}_{g}=\kappa_{g}\,\mathbf{J}_{m}\cdot\mathbf{B}_{g}
\]
\noindent This closes the entry by packaging the mass moment into the effective gravitomagnetic coupling coefficient.\par

\clearpage
\noindent\textbf{[ID 0325] Gravitomagnetic Coupling Term}\par
\vspace{0.35em}
\tzpsection{Dictionary}
\begingroup\small
\tzpsection{Definition}
Gravitomagnetic Coupling Term is a registry definition in Field Theory and Action Principles with secondary pressure from Manifold and Metric Dynamics.

% BEGIN TZP CONTEXTUAL MEANING
\tzpsection{Contextual Meaning}
\begingroup\footnotesize\setlength{\emergencystretch}{3em}\sloppy
\textbf{Formal statement:} This term represents the coupling between rotation and gravity, analogous to the magnetic field produced by moving charges. \textbf{Contextual meaning:} The entry reads Gravitomagnetic Coupling Term as a controlled technical headword whose equation layer, proof route, and registry role are interpreted together. \textbf{Derivable consequences:} The entry now states the gravitomagnetic coupling in explicit field-and-moment form instead of mixing it with stray code comments. \textbf{Lemma support:} Gravitomagnetic Field; Mass Integral; Coupling Strength.\par
\endgroup
% END TZP CONTEXTUAL MEANING

\tzpsection{Technical Interpretation}
The TZP topological charge is conserved because it is the flux of a divergence-free current through a hypersurface enclosing the zero-point structure. \% ========================================================= \% ========================================================= \% QUATERNARY DERIVATION LAYER -- ID 0326 TO ID 0330 \% Synthetic entries for absent identifiers.

\tzpsection{Equation Grounding}
Equation anchors: Configuration Space Structure **Statement**: System lives on manifold M = SO(3) times R to the power 3 times so(3)* times (R to the power 3)* **Logic Validation**: - Proposi,math sub methods ,validation; elastic sub energy = 0; and K * (d - d0)**2,elastic sub energy = 0. Proof route: show that the rotating mass distribution induces the stated gravitomagnetic field through the stated mass-moment integral. Formalization route: prove that the mass-moment integral defines the stated effective gravitomagnetic coupling coefficient.

\tzpsection{Interpretive Role}
The entry now states the gravitomagnetic coupling in explicit field-and-moment form instead of mixing it with stray code comments.
\endgroup

\tzpsection{TZP Thesaurus}
\begin{enumerate}[leftmargin=1.6em,itemsep=6pt,topsep=4pt]
\item \textbf{Mass-Current Field Generation}: The analogue to electromagnetism where circulating mass creates a dynamic field capable of deflecting passing matter.
\item \textbf{Ponderomotive Gravitational Coupling}: The specific constant defining the interaction strength between a localized spinning mass and ambient space.
\item \textbf{Kinetic-Gravitational Tension}: The stress field resulting from the interaction between local matter distribution velocities and global geometric curvature.
\end{enumerate}

\tzpsection{Encyclopedia Mathematica}
\begingroup\footnotesize\sloppy
The visual plate below is reserved for the source-specific equation and progression graphic for [ID 0325]. It is included only as a companion to the canonical equation, not as a substitute for the formal text.
\par\endgroup
\ifTZPDarkEdition
\tzpdictionaryvisualband{D:/TZPID/kdp_workflow/build/tikz_figures/ID0325_dictionary_figure_dark.pdf}
\fi

\clearpage
\begingroup\tzpsupportsetup
\noindent\textbf{\fontsize{10.8pt}{12.8pt}\selectfont [ID 0325] Constructive Logic and Proof Certificate}\par
\noindent\textbf{\fontsize{10pt}{12pt}\selectfont Gravitomagnetic Coupling Term}\par
\vspace{0.45em}
\begingroup
\footnotesize
\setlength{\parskip}{0.22em}
\setlength{\parindent}{0pt}
\noindent\textbf{Curry--Howard--Lambek Interpretation}\par
Under the Curry--Howard--Lambek correspondence, this registry entry is read as a typed proof
object: the stated source condition supplies the proposition, the derivation supplies the
morphism, and the proof certificate records the machine handles needed to check the obligation
in the formal registry.

\vspace{0.45em}
\noindent\textbf{CHL Proof}\par
\textbf{Statement:} this term represents the coupling between rotation and gravity, analogous to the
magnetic field produced by moving charges

\textbf{Logic Validation:}\\
\textbf{Proposition:} ``Gravitomagnetic Coupling Term is valid as a synchronization criterion under the coupling
conditions named by the source.''\\
\textbf{Proof:} The source statement describes coupled-oscillator behavior. The logical claim is
accepted only as a conditional synchronization result: coupling weights, frequency
spread, and order-parameter behavior determine whether coherence is increased, limited,
or dominated by a subset.

\textbf{VERDICT:} \ensuremath{\checkmark}\ \textbf{TRUE} --- conditional synchronization validation

\textbf{Formal Status:} \ensuremath{\checkmark}\ THEORETICAL -- source-backed CHL proof obligation

\textbf{Physical Status:} THEORETICAL -- requires model-specific empirical interpretation
\par\vspace{0.45em}
\noindent{\tiny\textbf{Isabelle/HOL.} \tzpplaincode{physics\_validity\_from\_model, external\_validity\_package, registry\_number\_matches, equation\_count\_matches}}\par
\ifTZPDarkEdition
\tzpproofvisualband{D:/TZPID/kdp_workflow/build/tikz_figures/ID0325_proof_certificate_figure_dark.pdf}
\fi
\endgroup
\endgroup
